\ifdefined\gkver\else\def\gkver{student}\fi
\documentclass[\gkver]{gaokaozhenti}
\begin{document}

\chapter{2026年1月浙江}

\begin{enumerate}
\item
%% number: 1
%% paperName: 2026年1月浙江选考第 1 题 3 分
%% typeId: 1
%% type: 单选题
%% chapter: 运动和力
%% point: 力学单位制
%% method: 无
%% score: 3
%% degree: 950
%% duplicateId: 0
%% body:
第六代移动通信技术（6G）使用的电磁波，部分处于太赫兹（THz）波段。$1\ \UTHZ = 10^{12}\ \UHz$，单位 THz 对应的物理量是 \xzanswer{C}

\fourchoices[answer=C]
{能量}
{功率}
{频率}
{波长}

%% answer: C
%% memo:
\memoanswer{
根据题干给出的单位关系 $1\ \UTHZ = 10^{12}\ \UHz$，其中 Hz（赫兹）是频率的单位，因此 THz（太赫兹）是频率的倍数单位，对应物理量为频率。

故选 C。
}

\item
%% number: 2
%% paperName: 2026年1月浙江选考第 2 题 3 分
%% typeId: 1
%% type: 单选题
%% chapter: 直线运动
%% point: 质点
%% method: 无
%% score: 3
%% degree: 950
%% duplicateId: 0
%% body:
下列问题中，图示物体可看成质点的是 \xzanswer{A}
\fourpicture[numstyle=chinese, h=2.2cm, labelA=20261月浙江02a, labelB=20261月浙江02b, labelC=20261月浙江02c, labelD=20261月浙江02d, capA=四川舰, capB=歼-35, capC=神舟二十二号, capD=蛟龙号]
{figs/02a.png}
{figs/02b.png}
{figs/02c.png}
{figs/02d.png}
\fourchoices[answer=A]
{研究图~\ref{20261月浙江02a} 中“四川舰”的航行路径}
{研究图~\ref{20261月浙江02b} 中“歼-35”战斗机的飞行姿态}
{研究图~\ref{20261月浙江02c} 中“神舟二十二号”载人飞船与空间站的对接方式}
{研究图~\ref{20261月浙江02d} 中“蛟龙号”潜水器完成任务出水后调整方位回舱过程}

%% answer: A
%% memo:
\memoanswer{
A．研究“四川舰”的航行路径时，“四川舰”的体积和形状可以忽略，故“四川舰”可以看成质点，故 A 正确；

B．研究“歼-35”战斗机的飞行姿态时，“歼-35”战斗机的体积和形状不可以忽略，故“歼-35”战斗机不可以看成质点，故 B 错误；

C．研究“神舟二十二号”载人飞船与空间站的对接方式时，“神舟二十二号”的体积和形状不可以忽略，故“神舟二十二号”不可以看成质点，故 C 错误；

D．研究“蛟龙号”潜水器完成任务出水后调整方位回舱过程时，“蛟龙号”的体积和形状不可以忽略，故“蛟龙号”不可以看成质点，故 D 错误。

故选 A。
}

\item
%% number: 3
%% paperName: 2026年1月浙江选考第 3 题 3 分
%% typeId: 1
%% type: 单选题
%% chapter: 运动和力
%% point: 牛顿第三定律
%% method: 无
%% score: 3
%% degree: 850
%% duplicateId: 0
%% body:
如图所示，钢架雪车运动员在具有阻力的倾斜赛道上滑行，则 \xzanswer{D}
\onepicture[w=5.2cm]{figs/03.jpg}
\fourchoices[answer=D]
{运动员在转弯时加速度为 0}
{运动员和钢架雪车整体机械能守恒}
{钢架雪车所受重力和赛道对钢架雪车的支持力是一对平衡力}
{钢架雪车对赛道的压力与赛道对钢架雪车的支持力是一对作用力和反作用力}

%% answer: D
%% memo:
\memoanswer{
A．运动员在转弯时一定有向心加速度，加速度不可能为零，故 A 错误；

B．倾斜赛道有阻力，阻力对运动员和钢架雪车做负功，运动员和钢架雪车整体机械能不守恒，故 B 错误；

C．钢架雪车所受重力竖直向下，赛道对钢架雪车的支持力垂直赛道向上，不在同一条直线上，不是一对平衡力，故 C 错误；

D．钢架雪车对赛道的压力与赛道对钢架雪车的支持力是一对作用力和反作用力，故 D 正确。

故选 D。
}

\item
%% number: 4
%% paperName: 2026年1月浙江选考第 4 题 3 分
%% typeId: 1
%% type: 单选题
%% chapter: 电场
%% point: 电场线
%% method: 无
%% score: 3
%% degree: 750
%% duplicateId: 0
%% body:
手机电容式触摸屏的核心部件可简化为平行板电容器。当手指靠近触摸屏时，电容器两极板和手指间的电场线分布如图所示。下列说法正确的是 \xzanswer{C}
\onepicture[w=6.1cm]{figs/04.png}
\fourchoices[answer=C]
{$A$ 点的电场强度大于 $B$ 点的电场强度}
{将一电子从 $A$ 点移到 $B$ 点，电子的电势能增大}
{极板上表面 $C$ 点的电势等于下表面 $D$ 点的电势}
{若电子在 $E$ 点释放，仅受静电力作用将沿电场线 $ab$ 运动}

%% answer: C
%% memo:
\memoanswer{
A．电场线越密的地方场强越大，电场线越疏的地方场强越小，由图可知，$A$ 点的电场强度小于 $B$ 点的电场强度，故 A 错误；

B．顺着电场线电势逐渐降低，由图可知 $A$ 点电势低于 $B$ 点电势，根据 $E_{p} = -e\varphi$ 可知将一电子从 $A$ 点移到 $B$ 点，电子的电势能减小，故 B 错误；

C．同一金属极板上，达到静电平衡后，电势处处相等，则极板上表面 $C$ 点的电势等于下表面 $D$ 点的电势，故 C 正确；

D．由于电场线 $ab$ 是曲线，则若电子在 $E$ 点释放，仅受静电力作用不可能沿电场线 $ab$ 运动，故 D 错误。

故选 C。
}

\item
%% number: 5
%% paperName: 2026年1月浙江选考第 5 题 3 分
%% typeId: 1
%% type: 单选题
%% chapter: 曲线运动
%% point: 人造地球卫星
%% method: 无
%% score: 3
%% degree: 650
%% duplicateId: 0
%% body:
已知行星的平均密度为 $\rho$，靠近行星表面运行的卫星做圆周运动的周期为 $T$。对于任何行星均为同一常量的是 \xzanswer{B}

\fourchoices[answer=B]
{$\rho T$}
{$\rho T^{2}$}
{$\rho^{2}T$}
{$\rho^{2}T^{3}$}

%% answer: B
%% memo:
\memoanswer{
根据万有引力定律和圆周运动规律，卫星在行星表面附近运行时，万有引力提供向心力 $G\frac{M}{R^{2}}=\frac{4\pi^{2}R}{T^{2}}$

行星平均密度 $\rho=\frac{M}{\frac{4}{3}\pi R^{3}}$

联立解得 $\rho=\frac{3\pi}{GT^{2}}$

A．$\rho T = \frac{3\pi}{GT}$，与 $T$ 有关，非常量，故 A 错误；

B．$\rho T^{2} = \frac{3\pi}{G}$，为常量，故 B 正确；

C．$\rho^{2}T = \frac{9\pi^{2}}{G^{2}T^{3}}$，与 $T$ 有关，非常量，故 C 错误；

D．$\rho^{2}T^{3} = \frac{9\pi^{2}}{G^{2}T}$，与 $T$ 有关，非常量，故 D 错误。

故选 B。
}

\item
%% number: 6
%% paperName: 2026年1月浙江选考第 6 题 3 分
%% typeId: 1
%% type: 单选题
%% chapter: 力物体的平衡
%% point: 三力平衡
%% method: 无
%% score: 3
%% degree: 650
%% duplicateId: 0
%% body:
如图所示，一盏重为 $G$ 的艺术灯用细绳悬挂，左右两侧细绳与水平方向夹角分别为 $45^{\circ}$ 和 $60^{\circ}$，细绳拉力分别为 $F_{1}$ 和 $F_{2}$。$A$ 和 $B$ 是左侧细绳两端点，$C$ 和 $D$ 分别是天花板和灯上的点，$CD$ 与 $AB$ 平行，则 \xzanswer{B}
\onepicture[w=6.5cm]{figs/06.png}
\fourchoices[answer=B]
{$F_{1}$ 大于 $F_{2}$}
{$F_{1}$ 和 $F_{2}$ 都小于 $G$}
{用细绳连接 $C$ 和 $B$ 后撤去 $AB$ 绳，可使灯位置不变}
{用细绳连接 $C$ 和 $D$ 后撤去 $AB$ 绳，可使灯位置不变}

%% answer: B
%% memo:
\memoanswer{
AB．对艺术灯受力分析，如图所示。

\onepicture[h=4.6cm]{figs/06_an1.png}

将 $F_{1}$、$F_{2}$ 和 $G$ 三力平移后，构成矢量三角形，如图。

\onepicture[h=4.0cm]{figs/06_an2.png}

根据矢量三角形可知 $F_{1} < F_{2} < G$，故 A 错误，B 正确；

CD．用细绳连接 $C$ 和 $B$ 后撤去 $AB$ 绳或用细绳连接 $C$ 和 $D$ 后撤去 $AB$ 绳，如图所示。

\onepicture[h=5.0cm]{figs/06_an3.png}

若艺术灯位置保持不变，则三力不能交汇一点，即无法保持平衡，故 CD 错误。

故选 B。
}

\item
%% number: 7
%% paperName: 2026年1月浙江选考第 7 题 3 分
%% typeId: 1
%% type: 单选题
%% chapter: 电场
%% point: 库仑定律
%% method: 无
%% score: 3
%% degree: 650
%% duplicateId: 0
%% body:
下列说法正确的是 \xzanswer{D}
\fourpicture[numstyle=chinese, h=3.2cm, labelA=20261月浙江07a, labelB=20261月浙江07b, labelC=20261月浙江07c, labelD=20261月浙江07d]
{figs/07a.png}
{figs/07b.png}
{figs/07c.png}
{figs/07d.png}
\fourchoices[answer=D]
{图~\ref{20261月浙江07a} 中线框在图示时刻的电流沿顺时针方向（俯视）}
{若图~\ref{20261月浙江07b} 中储罐内不导电液体液面上升，$LC$ 电路振荡周期减小}
{图~\ref{20261月浙江07c} 中单色光入射楔形透明膜时，频率越高，明暗条纹间距越大}
{用图~\ref{20261月浙江07d} 中扭秤探究电荷间相互作用力时，应使金属小球 $A$ 与 $C$ 带同种电荷}

%% answer: D
%% memo:
\memoanswer{
A．根据右手定则，图甲中线框在图示时刻的电流沿逆时针方向（俯视），A 错误；

B．若图乙中储罐内不导电液体液面上升，根据 $C=\frac{\varepsilon S}{4\pi kd}$，$\varepsilon$ 变大，则 $C$ 变大，根据 $T=2\pi\sqrt{LC}$ 可知，$LC$ 电路振荡周期增大，B 错误；

C．图丙中单色光入射楔形透明膜时，频率越高，波长越小，从透明膜前后表面反射的光的路程差相差一个波长的相邻两明纹的距离减小，即明暗条纹间距越小，C 错误；

D．用图丁中扭秤探究电荷间相互作用力时，应使金属小球 $A$ 与 $C$ 带同种电荷，两球由于排斥而使 $A$ 转动，D 正确。

故选 D。
}

\item
%% number: 8
%% paperName: 2026年1月浙江选考第 8 题 3 分
%% typeId: 1
%% type: 单选题
%% chapter: 电路
%% point: 电流强度
%% method: 无
%% score: 3
%% degree: 550
%% duplicateId: 0
%% body:
如图所示，一根带负电的塑料棒，长为 $l$、横截面积为 $S$、均匀分布有 $N$ 个电子（电子电荷量为 $-e$）。让棒垂直于匀强磁场 $B$，以速度 $v$ 沿轴向做匀速运动。下列说法正确的是 \xzanswer{B}
\onepicture[h=3.0cm]{\includesvg{figs/08}}
\fourchoices[answer=B]
{等效电流的方向与 $v$ 方向相同}
{等效电流的大小 $I = \frac{Nev}{l}$}
{棒产生的感应电动势 $E = Blv$}
{棒所受安培力的大小 $F = NeSvlB$}

%% answer: B
%% memo:
\memoanswer{
A．塑料棒带负电，则等效电流的方向与 $v$ 方向相反，A 错误；

B．等效电流的大小 $I = \frac{q}{\Delta t} = \frac{v\Delta t\cdot\frac{N}{l}e}{\Delta t} = \frac{veN}{l}$，B 正确；

C．棒产生的感应电动势 $E = B\cdot 2rv = 2Brv = 2Bv\sqrt{\frac{S}{\pi}}$（$r$ 为塑料棒的横截面的半径），C 错误；

D．棒所受安培力的大小 $F = BIl = NevB$，D 错误。

故选 B。
}

\item
%% number: 9
%% paperName: 2026年1月浙江选考第 9 题 3 分
%% typeId: 1
%% type: 单选题
%% chapter: 机械波
%% point: 波长频率波速
%% method: 无
%% score: 3
%% degree: 450
%% duplicateId: 0
%% body:
在地质探测中，可利用横波传播速度的不同，探测岩石密度信息。选择岩石分界面上的一点为原点、垂直分界面方向为 $x$ 轴，建立如图所示的坐标系。在坐标原点安装周期 $T = 1\times10^{-3}\ \Us$、振幅 $A = 2\Umm$ 的人工振源。$t = 0$ 时振源从平衡位置（$y = 0$）开始沿 $y$ 轴正方向振动，同时向两侧传播简谐横波。$t = 3\times10^{-3}\ \Us$ 时在岩石 Ⅰ（$x < 0$）中的波恰好到达 $x = -3\Um$ 处，岩石 Ⅱ（$x > 0$）中的波速为 $v_{\text{Ⅱ}} = 3\times10^{3}\Ums$，则 \xzanswer{C}
\onepicture[w=9.0cm]{figs/09.png}
\fourchoices[answer=C]
{岩石 Ⅰ 和岩石 Ⅱ 中两波波长之比为 $3:1$}
{$t = 1.5\times10^{-3}\ \Us$ 时，$x = -1\Um$ 处质点的振动方向沿 $y$ 轴正方向}
{在 $0\sim3\times10^{-3}\ \Us$ 内，$x = 1.25\Um$ 处质点经过的路程为 $21\Umm$}
{增大振源的振动周期，岩石 Ⅱ 中的波速将变小}

%% answer: C
%% memo:
\memoanswer{
A．岩石 Ⅰ 中，横波传播速度 $v_{1}=\frac{\left|\Delta x\right|}{\Delta t}=1\times10^{3}\Ums$

波长公式 $\lambda = vT$，两波周期相同，故波长比等于波速比 $\frac{\lambda_{1}}{\lambda_{\text{Ⅱ}}} = \frac{v_{1}}{v_{\text{Ⅱ}}} = \frac{1\times10^{3}}{3\times10^{3}} = \frac{1}{3}$，A 错误；

B．$t = 0$ 时振源从平衡位置（$y = 0$）开始沿 $y$ 轴正方向振动，波传到 $x = -1\Um$ 的时间 $t_{1} = \frac{1\Um}{1\times10^{3}\Ums} = 1\times10^{-3}\Us$

在 $t = 1.5\times10^{-3}\Us$ 时，可知 $x = -1\Um$ 处质点的已经振动的时间 $\Delta t = 0.5\times10^{-3}\Us = \frac{T}{2}$

振源起振方向为 $y$ 轴正方向，经过半个周期后，质点振动方向为 $y$ 轴负方向，B 错误；

C．波传到 $x = 1.25\Um$ 的时间 $t_{2} = \frac{1.25\Um}{3\times10^{3}\Ums} = \frac{5}{12}\times10^{-3}\Us$

该质点振动的总时间 $\Delta t' = \left(3 - \frac{5}{12}\right)\times10^{-3}\Us = \frac{31}{12}\times10^{-3}\Us = \frac{31}{12}T$

即 2 个全振动和剩余时间 $\frac{7}{12}T$，质点从平衡位置（$y=0$）开始向 $y$ 轴正方向振动，则该质点经过的路程为 $s = 2.5\times4A + A\left|\sin\left(2\pi\times\frac{7}{12}\right)\right| = 21\Umm$，C 正确；

D．机械波在介质中的传播速度由介质本身的性质（如密度、弹性模量）决定，与振源的周期、频率无关，故波速不变，D 错误。

故选 C。
}

\item
%% number: 10
%% paperName: 2026年1月浙江选考第 10 题 3 分
%% typeId: 1
%% type: 单选题
%% chapter: 电磁感应
%% point: 电磁感应中图象
%% method: 无
%% score: 3
%% degree: 450
%% duplicateId: 0
%% body:
如图~\ref{20261月浙江10a} 所示，半径为 $r_{1}$、横截面半径为 $r_{2}$（$r_{2}\ll r_{1}$）、匝数为 $N$ 的圆环形螺线管通有电流 $I$，管内产生磁感应强度 $B = aI$（$a$ 为常量）的匀强磁场。管外磁场近似为 0，小明用电阻为 $R$ 的一段漆包线缠绕螺线管一圈后，并成双股线再缠绕螺线管两圈，最后将两端头短接，形成特殊线圈 A。若电流 $I$ 随时间 $t$ 变化的关系如图~\ref{20261月浙江10b} 所示，则 \xzanswer{D}
\twopicture[labelA=20261月浙江10a, labelB=20261月浙江10b, h=3.9cm]
{figs/10a.png}
{figs/10b.png}
\fourchoices[answer=D]
{$t = 2t_{0}$ 时，螺线管的自感电动势 $E = \frac{N\pi r_{1}^{2}aI_{0}}{t_{0}}$}
{$t = 5t_{0}$ 时，线圈 A 的感应电动势 $E = \frac{3\pi r_{2}^{2}aI_{0}}{t_{0}}$}
{在 $0\sim2t_{0}$ 内，通过线圈 A 的电荷量 $q = \frac{N\pi r_{2}^{2}aI_{0}}{R}$}
{在 $0\sim5t_{0}$ 内，线圈 A 产生的焦耳热 $Q = \frac{3(\pi r_{2}^{2}aI_{0})^{2}}{Rt_{0}}$}

%% answer: D
%% memo:
\memoanswer{
A．螺线管的自感电动势公式为 $E = N\frac{\Delta\Phi}{\Delta t} = NS\left|\frac{\Delta B}{\Delta t}\right|$

其中，螺线管的横截面积 $S = \pi r_{2}^{2}$，磁感应强度 $B = aI$。在 $t_{0}\sim2t_{0}$ 时间段内，电流变化率为 $\frac{\Delta I}{\Delta t} = \frac{-I_{0}-I_{0}}{2t_{0}-t_{0}} = -\frac{2I_{0}}{t_{0}}$

磁感应强度变化率为 $\frac{\Delta B}{\Delta t} = a\frac{\Delta I}{\Delta t} = -\frac{2aI_{0}}{t_{0}}$

代入自感电动势公式 $E = \frac{2N\pi r_{2}^{2}aI_{0}}{t_{0}}$，A 错误；

B．线圈 A 的有效匝数是 1，在 $4t_{0}\sim6t_{0}$ 时间段内，电流变化率为 $\frac{\Delta I}{\Delta t} = \frac{2I_{0}}{2t_{0}} = \frac{I_{0}}{t_{0}}$

线圈 A 的感应电动势 $E = nS\frac{\Delta B}{\Delta t} = \frac{\pi r_{2}^{2}aI_{0}}{t_{0}}$，B 错误；

C．流过线圈 A 的电荷量 $q = \overline{I}\Delta t = \frac{\overline{E}}{R}\Delta t = \frac{\Delta\Phi_{\mathrm{A}}}{R}$，其中 $\Delta\Phi_{\mathrm{A}}$ 是穿过线圈 A 的总磁通量变化量。在 $0\sim2t_{0}$ 内，螺线管电流恒定，穿过线圈 A 的总磁通量不变，通过线圈 A 的电荷量为 0。在 $t_{0}\sim2t_{0}$ 内，电流从 $I_{0}$ 变为 0，所以 $\Delta I = 0 - I_{0} = -I_{0}$。

在 $0\sim2t_{0}$ 内，通过线圈 A 的电荷量 $q = \frac{\Delta\Phi_{\mathrm{A}}}{R} = \frac{\pi r_{2}^{2}aI_{0}}{R}$，C 错误；

D．在 $0\sim5t_{0}$ 内，线圈 A 产生焦耳热只有在电流变化的区间才产生感应电动势和焦耳热，即 $t_{0}\sim3t_{0}$ 和 $4t_{0}\sim5t_{0}$ 时间段。$t_{0}\sim3t_{0}$ 时长为 $2t_{0}$，线圈 A 的感应电动势 $E_{1}=nS\frac{\Delta B}{\Delta t}=\frac{\pi r_{2}^{2}aI_{0}}{t_{0}}$

产生的焦耳热 $Q_{1}=\frac{E_{1}^{2}}{R}\cdot2t_{0}=2\frac{\left(\pi r_{2}^{2}aI_{0}\right)^{2}}{Rt_{0}}$

在 $4t_{0}\sim5t_{0}$ 时间段内，感应电动势 $E_{2}=nS\frac{\Delta B}{\Delta t}=\frac{\pi r_{2}^{2}aI_{0}}{t_{0}}$

产生的焦耳热 $Q_{2}=\frac{E_{2}^{2}}{R}t_{0}=\frac{\left(\pi r_{2}^{2}aI_{0}\right)^{2}}{Rt_{0}}$

在 $0\sim5t_{0}$ 内，线圈 A 产生的焦耳热 $Q = Q_{1} + Q_{2} = \frac{3(\pi r_{2}^{2}aI_{0})^{2}}{Rt_{0}}$，D 正确。

故选 D。
}

\item
%% number: 11
%% paperName: 2026年1月浙江选考第 11 题 4 分
%% typeId: 2
%% type: 多选题
%% chapter: 曲线运动
%% point: 竖直平面圆周运动
%% method: 无
%% score: 4
%% degree: 750
%% duplicateId: 0
%% body:
下列说法正确的是 \xzanswer{AB}

\fourchoices[answer=AB]
{变化的电场会产生磁场}
{汽车经过凹形路面最低点时处于超重状态}
{相对论时空观认为运动物体的长度与速度无关}
{秦山核电站通过提取海水中的氘和氚进行核反应获取能量}

%% answer: AB
%% memo:
\memoanswer{
A．根据麦克斯韦电磁理论可知，变化的电场会产生磁场，故 A 正确；

B．汽车经过凹形路面最低点时，其加速度指向圆心向上，可知此时汽车处于超重状态，故 B 正确；

C．相对论时空观认为运动物体的长度与速度有关，物体的长度会随着物体的速度不同而改变，故 C 错误；

D．通过提取海水中的氘和氚进行核反应属于轻核聚变，而秦山核电站利用重核裂变获取能量，故 D 错误。

故选 AB。
}

\item
%% number: 12
%% paperName: 2026年1月浙江选考第 12 题 4 分
%% typeId: 2
%% type: 多选题
%% chapter: 光学
%% point: 全反射
%% method: 无
%% score: 4
%% degree: 450
%% duplicateId: 0
%% body:
如图~\ref{20261月浙江12a} 所示，三块同质有机玻璃板甲、乙、丙，拼接形成两层同心圆弧空气膜。现让一束激光从 $P$ 点射入，并在空气膜 $Q$ 处注入一滴油，呈现图~\ref{20261月浙江12b} 所示的光路（其中 $\alpha$、$\beta$ 为相应光线的入射角）。已知空气折射率为 1，光在空气中的传播速度为 $c$。下列说法正确的是 \xzanswer{AC}
\twopicture[labelA=20261月浙江12a, labelB=20261月浙江12b, h=3.4cm]
{figs/12a.png}
{figs/12b.png}
\fourchoices[answer=AC]
{有机玻璃的折射率 $n \geq \frac{1}{\sin\beta}$}
{光在有机玻璃中的传播速度 $v \geq c\sin\alpha$}
{有机玻璃的折射率可能大于油的折射率}
{若在 $R$ 处注入同种油滴，光线不会从 $R$ 处进入乙}

%% answer: AC
%% memo:
\memoanswer{
A．光线在外层空气膜发生全反射，则 $\sin\beta \geq \sin C = \frac{1}{n}$，即有机玻璃的折射率 $n \geq \frac{1}{\sin\beta}$，A 正确；

B．光线在内层空气膜发生全反射，则 $\sin\alpha \geq \sin C = \frac{1}{n_{\text{玻璃}}}$，$n_{\text{玻璃}} = \frac{c}{v}$，可得光在有机玻璃中的传播速度 $v \leq c\sin\alpha$，B 错误；

CD．在空气膜 $Q$ 处注入一滴油后，光线能通过油膜界面，不再发生全反射，则此时临界角 $\sin\alpha \leq \sin C' = \frac{n_{\text{油}}}{n_{\text{玻璃}}}$，则此时有机玻璃的折射率大于油的折射率，若在 $R$ 处注入同种油滴，光线同样可沿直线通过油膜界面，光线会从 $R$ 处进入乙，C 正确，D 错误。

故选 AC。
}

\item
%% number: 13
%% paperName: 2026年1月浙江选考第 13 题 4 分
%% typeId: 1
%% type: 单选题
%% chapter: 磁场
%% point: 带电粒子在磁场中的运动
%% method: 无
%% score: 4
%% degree: 350
%% duplicateId: 0
%% body:
在充满液态氢的气泡室中存在方向垂直图示平面、磁感应强度为 $B$ 的匀强磁场。一束 $\gamma$ 射线通过气泡室，其中一个 $\gamma$ 光子将一个氢原子打出一个电子（$\mathrm{e}^{-}$），同时自身转变为一对正负电子对（$\mathrm{e}^{+}-\mathrm{e}^{-}$）。三个电子在气泡室中的径迹如图所示（氢原子和产生的质子均可视为静止），开始时，三条径迹共切于 $O$ 点，其半径分别为 $r_{1}$、$r_{2}$ 和 $r_{3}$。假设电子在气泡室中所受阻力大小正比于速率，比例系数为 $k$，方向与速度方向相反。沿径迹 1、2、3 运动的电子速度减为 0 时的位置分别位于 $M$（图中未标出）、$P$、$Q$ 三点。已知电子质量为 $m_{e}$，元电荷为 $e$，光速为 $c$。下列说法正确的是 \xzanswer{BC}
\onepicture[h=4.6cm]{\includesvg{figs/13}}
\fourchoices[answer=BC]
{$\gamma$ 光子的能量小于 $2m_{e}c^{2}$}
{$\gamma$ 光子的动量大小为 $eB(r_{1} + r_{2} + r_{3})$}
{$O$ 与 $P$、$Q$ 的距离之比 $OP:OQ = r_{2}:r_{3}$}
{沿径迹 1 运动的电子总路程为 $\frac{\sqrt{(eB)^{2} + k^{2}}}{k}r_{1}$}

%% answer: BC
%% memo:
\memoanswer{
A．因 $\gamma$ 光子转化为一对正负电子对应的能量为 $2m_{e}c^{2}$，同时 $\gamma$ 光子将一部分能量传递给被打出的电子，可知 $\gamma$ 光子的能量一定大于 $2m_{e}c^{2}$，A 错误；

B．电子在磁场中做圆周运动，根据 $evB = m\frac{v^{2}}{r}$

可得 $p = mv = eBr$

则沿轨迹 1、2、3 运动的电子对应的动量分别为 $p_{1}=eBr_{1}$、$p_{2}=eBr_{2}$、$p_{3}=eBr_{3}$

由动量守恒可知 $p_{\gamma}=p_{1}+p_{2}+p_{3}$

可得 $\gamma$ 光子的动量大小为 $p_{\gamma}=eB(r_{1}+r_{2}+r_{3})$，B 正确；

C．正负电子在磁场运动时受阻力作用速度逐渐减小，则做圆周运动的半径逐渐减小，即电子做螺旋运动，将电子的轨迹分成无数小段，每一小段均可看做是一段圆弧，因速度方向总是沿对应的一小段圆弧的切线方向，电子在每一小段上做圆周运动的圆心是固定的（该位置为电子刚开始做圆周运动时的圆心位置，对沿 2、3 轨迹的电子分别为 $P$ 点和 $Q$ 点），电子最终停止运动时将停止在该圆心位置，可知 $O$ 与 $P$、$Q$ 的距离之比 $OP:OQ = r_{2}:r_{3}$，C 正确；

D．沿径迹 1 运动的电子初速度为 $v_{1} = \frac{eBr_{1}}{m}$

电子受的阻力 $f = kv$

由动量定理 $-kv\Delta t = 0 - m\Delta v$

求和可得 $\sum kv\Delta t = \sum m\Delta v$

即 $ks = mv_{1}$

总路程为 $s = \frac{eB}{k}r_{1}$，D 错误。

故选 BC。
}

\item
%% number: 14
%% paperName: 2026年1月浙江选考第 14 题 2 分
%% typeId: 4
%% type: 实验题
%% chapter: 气体性质
%% point: 用单分子油膜估测分子大小
%% method: 无
%% score: 2
%% degree: 650
%% duplicateId: 0
%% body:
实验题（Ⅰ、Ⅱ、Ⅲ 三题共 14 分）

Ⅰ．（6 分）“在探究平抛运动的特点”实验中，

\begin{enumerate}
	\item
	为探究水平方向分运动特点，应选用图~\ref{20261月浙江14a} 中的 \tkanswer[1.2cm]{乙} （选填“甲”或“乙”）装置；
	\item
	采用图~\ref{20261月浙江14b} 所示装置进行实验。将一张白纸和复写纸固定在装置的背板上，钢球落到倾斜的挡板后挤压复写纸，在白纸上留下印迹。下列说法正确的是 \tkanswer[1.2cm]{AB} 。

	A．调节装置使其背板竖直

	B．调节斜槽使其末端切线水平

	C．以斜槽的末端在白纸上的投影点为坐标原点

	D．钢球在斜槽静止释放的高度应等间距下降

	\item
	如图~\ref{20261月浙江14c} 所示，将实验中记录的印迹用平滑曲线连接，其中抛出点为坐标原点，$A$ 点（$11.0\Ucm$，$15.8\Ucm$）是记录的印迹，$B$ 点（$11.8\Ucm$，$19.6\Ucm$）是曲线上的一个点，为得到小球的水平速度，应取 \tkanswer[1.2cm]{B} （选填“$A$”或“$B$”）点进行计算，可得水平速度 $v_{0}=$ \tkanswer[1.4cm]{0.59} $\Ums$。（$g$ 取 $9.8\Ums^{2}$，所得结果保留两位有效数字）
\end{enumerate}
\threepicture[labelA=20261月浙江14a, labelB=20261月浙江14b, labelC=20261月浙江14c, numA=1, numB=2, numC=3, hA=3.4cm, hB=3.4cm, hC=5.2cm, align=right]
{figs/14a.png}
{figs/14b.png}
{figs/14c.png}
Ⅱ．（6 分）在“探究变压器原、副线圈电压与匝数的关系”实验中，

\begin{enumerate}
	\item
	小明先用多用电表测量图~\ref{20261月浙江14d} 变压器的“0”、“1 400”接线柱间的电阻。选择开关位置如图~\ref{20261月浙江14e}，经规范操作，指针位置如图~\ref{20261月浙江14f}，测得阻值为 \tkanswer[1.4cm]{34.0} $\UO$。
	\item
	用匝数 $n_{1} = 200$ 和 $n_{2} = 400$ 的两线圈进行实验，分别测得两端电压为 $U_{1}$ 和 $U_{2}$，记录于下表。下列说法正确的是 \tkanswer[1.2cm]{BD} （多选）；

	\begin{center}
	\begin{tabular}{|c|c|c|c|c|c|}
	\hline
	$U_{1}/\UV$ & 1.02 & 2.20 & 3.24 & 4.28 & 5.36 \\ \hline
	$U_{2}/\UV$ & 2.12 & 4.52 & 6.64 & 8.78 & 11.02 \\ \hline
	\end{tabular}
	\end{center}

	A．与 $n_{1}$ 对应的是副线圈

	B．与 $n_{2}$ 对应的是副线圈

	C．实验中采用低压直流电源

	D．多用电表选择开关应调至交流电压挡

	\item
	小明将两个线圈按图~\ref{20261月浙江14g} 方式上下叠放。下层线圈输入电压信号如图~\ref{20261月浙江14h} 所示，上层线圈与示波器相连，则示波器上显示的波形为 \tkanswer[1.2cm]{C} 。
\end{enumerate}
\threepicture[labelA=20261月浙江14d, labelB=20261月浙江14e, labelC=20261月浙江14f, numA=1, numB=2, numC=3, h=3.0cm]
{figs/14d.png}
{figs/14e.png}
{figs/14f.png}
\twopicture[topsep=8pt, labelA=20261月浙江14g, labelB=20261月浙江14h, numA=4, numB=5, h=2.6cm]
{figs/14g.png}
{figs/14h.png}
\threepicture[topsep=8pt, labelprefix={}, numA=A, numB=B, numC=C, h=2.2cm]
{figs/14i.png}
{figs/14j.png}
{figs/14k.png}
Ⅲ．（2 分）在“油膜法估测油酸分子的大小”实验中，

\begin{enumerate}
	\item
	以下说法正确的是 \tkanswer[1.2cm]{B}

	A．图示油膜形状是由于撒粉太少引起的

	B．按图示油膜面积进行计算，测得油酸分子直径偏大

	C．油酸酒精溶液放置长时间后使用，测得油酸分子直径偏大

	\item
	测得一滴油酸酒精溶液所含纯油酸的体积为 $V$，根据画有油膜轮廓的玻璃板上的坐标方格，数出轮廓范围内正方形的个数，整格的为 $N_{1}$ 个，多于半格不足整格的数量为 $N_{2}$ 个，已知每格的面积为 $S$，则油酸分子的直径为 \tkanswer[3.2cm]{$\frac{V}{(N_{1}+N_{2})S}$} 。
\end{enumerate}
\onepicture[align=right, w=4.2cm]{figs/14l.jpg}
%% answer:
\jdanswer{
\begin{enumerate}
	\item
	乙
	\item
	AB
	\item
	B，$0.59\Ums$
\end{enumerate}
\begin{enumerate}
	\item
	$34.0\UO$
	\item
	BD
	\item
	C
\end{enumerate}
\begin{enumerate}
	\item
	B
	\item
	$\frac{V}{(N_{1}+N_{2})S}$
\end{enumerate}
}
%% memo:
\memoanswer{
（1）为研究水平方向分运动特点，需要将平抛运动与匀速直线运动进行对比。故应该选装置乙。

（2）A．调节装置使其背板竖直，才能保证小球落在背板上的痕迹准确反映平抛轨迹，A 正确。

B．调节斜槽末端切线水平，才能保证小球抛出时初速度水平，B 正确。

C．坐标原点应选小球在斜槽末端时球心的位置，不是斜槽末端的投影点，C 错误。

D．每次释放小球的高度应相同，保证初速度一致，不需要等间距下降，D 错误。

故选 AB。

（3）由图可知因为 $B$ 点更接近曲线轨迹上，所以为得到小球的水平速度应该选 $B$ 点进行计算；

从原点到 $B$ 点，小球竖直方向上做自由落体运动，时间为 $t=\sqrt{\frac{2y_{B}}{g}}=\sqrt{\frac{2\times0.196}{9.8}}=0.2\Us$

所以 $v_{0}=\frac{x_{B}}{t}=\frac{0.118}{0.2}=0.59\Ums$。

Ⅱ．（1）从图 2 可知，多用电表选择开关置于电阻$\times$10 挡位，图 3 中指针指向 3.4，所以测得阻值为 $R=3.4\times10=34\UO$。

（2）AB．根据变压器电压与匝数的关系，匝数多的线圈电压高；而表格中 $U_{2} > U_{1}$，所以 $n_{2}$ 对应副线圈，故 B 正确，A 错误。

C．变压器的工作原理是电磁感应，需要交流电源，不能用直流电源，C 错误。

D．变压器输出的是交流电，测量其电压时，多用电表选择开关应调至交流电压挡，D 正确。

故选 BD。

（3）由题意可知下层线圈输入的是锯齿波电压，它的变化率（斜率）在每个周期内是恒定的；根据电磁感应原理可知上层线圈的感应电动势与磁通量的变化率成正比，也就是与输入电压的变化率成正比；输入电压的斜率恒定，所以感应电动势的大小也恒定，对应的波形是方波。故选 C。

Ⅲ．（1）A．撒粉太少会导致油膜扩散时没有足够阻力，可能扩散得过大甚至溢出，或者边界模糊不清；而图中油膜呈现“尖刺”、“星芒”状，通常是撒粉太多或不均匀，导致油膜在局部被阻挡，无法均匀铺展，从而形成这种不规则形状，故 A 错误；

B．按图示油膜面积计算时，测量的面积比实际油膜的面积会更小，根据 $d=\frac{V}{S}$ 可知测得油酸分子直径偏大，故 B 正确；

C．油酸酒精溶液放置长时间后使用，因为酒精的挥发，在溶液中的油酸浓度会偏大，油酸的体积会偏大，但实际计算时仍按原来的浓度计算，则计算所用的油酸体积偏小，根据 $d=\frac{V}{S}$ 可知测得油酸分子直径偏小，故 C 错误。

故选 B。

（2）多于半格不足整格的按整格计算，故总的格子数为 $N_{1} + N_{2}$，则总面积为 $S_{\text{总}} = S(N_{1} + N_{2})$，则油酸分子直径为 $d = \frac{V}{S_{\text{总}}} = \frac{V}{(N_{1} + N_{2})S}$。
}

\item
%% number: 15
%% paperName: 2026年1月浙江选考第 15 题 8 分
%% typeId: 6
%% type: 计算题
%% chapter: 气体性质
%% point: 玻意耳定律
%% method: 无
%% score: 8
%% degree: 650
%% duplicateId: 0
%% body:
如图所示，导热良好的瓶内，用一质量为 $m_{1}$、横截面积为 $S$ 的活塞封闭一定质量的理想气体，活塞能无摩擦滑动，在活塞上方有质量为 $m_{2}$ 的液体。初始时，瓶内气体处于状态 $A$，体积为 $V_{A}$。将一根质量不计的细管插入液体，液体在细管中上升到一定高度后保持静止，随后通过细管缓慢吸走全部液体，此时瓶内气体处于状态 $B$。环境温度保持不变，从状态 $A$ 到状态 $B$ 过程中，气体吸收热量 $Q = 2.05\UJ$。已知 $m_{1} = 0.4\Ukg$，$m_{2} = 2\Ukg$，$S = 40\Ucmq$，$V_{A} = 400\Ucmc$，大气压强 $p_{0} = 0.99\times10^{5}\UPa$，$g = 10\Ums^{2}$。
\begin{enumerate}
	\item
	图中液体 \tkanswer[1.2cm]{浸润} （选填“浸润”或“不浸润”）管壁，若细管仅内径变小，与原细管相比，管内液面将 \tkanswer[1.2cm]{升高} （选填“升高”、“不变”或“降低”）；
	\item
	求气体在状态 $B$ 时的体积 $V_{B}$；
	\item
	求气体从状态 $A$ 到状态 $B$ 过程中对外做的功 $W_{AB}$。
\end{enumerate}
\onepicture[align=right, w=3.6cm]{figs/15.png}
%% answer: 
\jdanswer{
\begin{enumerate}
	\item
	浸润，升高

	\item
	$420\Ucmc$

	\item
	$2.05\UJ$

\end{enumerate}
}
%% memo:
\memoanswer{
（1）图中管中液面上升且液面呈现凹状，则液体浸润管壁，若细管仅内径变小，与原细管相比，毛细现象更加明显，管内液面升高。

（2）初态对活塞以及上面的液体分析可知气体压强
\[ p_{1}=p_{0}+\frac{(m_{1}+m_{2})g}{S}=1.05\times10^{5}\UPa \]
末态吸走液体后气体的压强为
\[ p_{2}=p_{0}+\frac{m_{1}g}{S}=1.00\times10^{5}\UPa \]
根据玻意耳定律可知 $p_{1}V_{A}=p_{2}V_{B}$

解得气体在状态 $B$ 时的体积为
\[ V_{B}=\frac{p_{1}V_{A}}{p_{2}}=\frac{1.05\times10^{5}\times400}{1.00\times10^{5}}\Ucmc=420\Ucmc \]

（3）气体从状态 $A$ 到状态 $B$ 过程中气体温度不变，则根据热力学第一定律 $\Delta U = W + Q$，其中 $\Delta U = 0$。

$W = -Q = -2.05\UJ$

即气体对外做的功 $2.05\UJ$。
}

\item
%% number: 16
%% paperName: 2026年1月浙江选考第 16 题 11 分
%% typeId: 6
%% type: 计算题
%% chapter: 电路
%% point: 传感器
%% method: 无
%% score: 11
%% degree: 450
%% duplicateId: 0
%% body:
测量局域磁场，科学家基于电阻应变片开发出一种磁场检测芯片，其简化结构如图~\ref{20261月浙江16a} 所示。长度均为 $l$、通有恒定电流 $I_{0}$（方向相反）的两刚性金属杆 $ab$、$cd$，与具有良好弹性的绝缘悬梁 $OA$、$OD$ 构成“$H$”形支架，对称固定于底座 $O$ 处。在悬梁上、下表面对称安装四个相同的电阻应变片（各自引出两导线），其阻值分别为 $R_{1}$、$R_{2}$、$R_{3}$ 和 $R_{4}$，将它们按图~\ref{20261月浙江16b} 方式与电动势为 $E$ 的电源（不计内阻）相连。未加磁场时，支架处于水平平衡状态，此时 $R_{1} = R_{2} = R_{3} = R_{4} = R_{0}$，测得 $e$、$f$ 两端的电势差为 0。现施加待测磁场，其方向水平向右、且垂直于金属杆，则金属杆 $ab$、$cd$ 受安培力作用，使悬梁 $OA$、$OD$ 产生形变，四个应变片的阻值发生相应变化，其变化量的绝对值均为 $\Delta R$，此时测得 $e$、$f$ 两端的电势差为 $U_{ef}$，从而得到待测磁场磁感应强度 $B$ 的大小。
\begin{enumerate}
	\item
	判断金属杆 $ab$ 和 $cd$ 所受安培力的方向；
	\item
	写出上述四个电阻的阻值（用 $R_{0}$ 和 $\Delta R$ 表示）；
	\item
	已知电阻变化量和所受的安培力成正比关系，且比例系数为 $\alpha$，求 $U_{ef}$ 与 $B$ 之间的关系。
\end{enumerate}
\twopicture[labelA=20261月浙江16a, labelB=20261月浙江16b, h=3.4cm, align=right]
{figs/16a.png}
{figs/16b.png}
%% answer: 
\jdanswer{
\begin{enumerate}
	\item
	$ab$ 竖直向下，$cd$ 竖直向上

	\item
	$R_{1} = R_{0} + \Delta R$，$R_{2} = R_{0} - \Delta R$，$R_{3} = R_{0} - \Delta R$，$R_{4} = R_{0} + \Delta R$

	\item
	$B = \frac{R_{0}U_{ef}}{\alpha EI_{0}l}$

\end{enumerate}
}
%% memo:
\memoanswer{
（1）根据左手定则可知 $ab$ 所受安培力方向竖直向下；$cd$ 所受安培力方向竖直向上。

（2）由题意可知 $ab$ 向上弯曲，使 $R_{1}$ 被拉伸（阻值增大）、$R_{3}$ 被压缩（阻值减小），故 $R_{1}=R_{0}+\Delta R$，$R_{3}=R_{0}-\Delta R$

$cd$ 向下弯曲，使 $R_{2}$ 被压缩（阻值减小）、$R_{4}$ 被拉伸（阻值增大），故 $R_{2}=R_{0}-\Delta R$，$R_{4}=R_{0}+\Delta R$

（3）由图可知 $R_{1}$ 与 $R_{2}$ 串联，$R_{3}$ 与 $R_{4}$ 串联，两条支路并联。

上支路总电阻为 $R_{\text{上}}=R_{1}+R_{2}=(R_{0}+\Delta R)+(R_{0}-\Delta R)=2R_{0}$

电流 $I_{\text{上}}=\frac{E}{2R_{0}}$

$f$ 点电势 $\varphi_{f}=I_{\text{上}}R_{2}=\frac{E(R_{0}-\Delta R)}{2R_{0}}$

下支路总电阻为 $R_{\text{下}}=R_{3}+R_{4}=(R_{0}-\Delta R)+(R_{0}+\Delta R)=2R_{0}$

电流 $I_{\text{下}}=\frac{E}{2R_{0}}$

$e$ 点电势 $\varphi_{e}=I_{\text{下}}R_{4}=\frac{E(R_{0}+\Delta R)}{2R_{0}}$

$e$、$f$ 两点间的电势差绝对值 $U_{ef}=\varphi_{e}-\varphi_{f}=\frac{E\Delta R}{R_{0}}$

安培力与 $\Delta R$ 的关系：$\Delta R = \alpha F$，而 $F = BI_{0}l$，所以 $\Delta R = \alpha BI_{0}l$

联立得 $B = \frac{R_{0}U_{ef}}{\alpha EI_{0}l}$
}

\item
%% number: 17
%% paperName: 2026年1月浙江选考第 17 题 12 分
%% typeId: 6
%% type: 计算题
%% chapter: 机械振动
%% point: 等效单摆
%% method: 无
%% score: 12
%% degree: 350
%% duplicateId: 0
%% body:
如图所示。一宽度为 $d$ 的光滑长方形平板 $MNOP$，长边 $MN$、$PQ$ 分别平滑连接半径均为 $r$ 的光滑圆弧面，形成“U”形槽，将其整体固定在水平地面上。现有质量为 $m$ 的物块 a，从圆弧面上相对平板竖直高度为 $h$ 的 $A$ 点静止下滑（$h \ll r$），途经圆弧面上最低点 $B$，平板上有一质量为 $\frac{m}{3}$ 的物块 b 与 $MN$ 成 $45^{\circ}$ 角从 $O$ 点滑入圆弧面，第一次到达最高点时恰好与同时到达最高点的物块 a 发生弹性碰撞。两物块均为质点。
\begin{enumerate}
	\item
	求物块 a 第一次经过 $B$ 点时速度大小 $v_{0}$ 和所受支持力大小 $F_{N}$；
	\item
	从 $A$ 到 $B$ 的过程：物块 a 相对于 $B$ 点位移为 $x$，求其所受回复力 $F$ 与 $x$ 的关系式；
	\item
	求物块 b 的初速度大小 $v_{b}$ 以及碰撞后瞬间物块 a 的速度大小 $v_{a}$；
	\item
	若 $h = 0.032\Um$，$r = 10\Um$，$d = 0.4\Um$，要使物块 a 从 $NQ$ 之间滑离，求 $BQ$ 间距 $L$ 的范围。
\end{enumerate}
\onepicture[align=right, w=7.0cm]{figs/17.png}
%% answer: 
\jdanswer{
\begin{enumerate}
	\item
	$v_{0} = \sqrt{2gh}$，$F_{N} = m(g + \frac{2gh}{r})$，方向竖直向上

	\item
	$F = -\frac{mg}{r}x$

	\item
	$v_{b} = 2\sqrt{gh}$，$v_{a} = \sqrt{\frac{gh}{2}}$

	\item
	$[0.2\pi + n(0.4\pi + 0.2)]\Um < L < [0.2\pi + 0.2 + n(0.4\pi + 0.2)]\Um$，（$n = 0,1,2,3,4\cdots\cdots$）

\end{enumerate}
}
%% memo:
\memoanswer{
（1）对 a 物块下滑过程中根据动能定理有 $mgh = \frac{1}{2}mv_{0}^{2}$

可得 $v_{0} = \sqrt{2gh}$；

在 B 点根据牛顿第二定律有 $F_{N}-mg=m\frac{v_{0}^{2}}{r}$

可得 $F_{N}=m(g+\frac{2gh}{r})$，方向竖直向上；

（2）如图。由于 $h \ll r$，滑块 a 所受回复力 $F$，则 $F = -mg\sin\theta = -\frac{mg}{r}x$

可知滑块受到的回复力 $F$ 与 $x$ 成正比，方向与 $x$ 相反，因此滑块 a 从释放到第一次到达最低点的运动是简谐运动。

\onepicture[h=5.0cm]{figs/17_an1.png}

（3）滑块 b 在圆弧形斜面上垂直槽轴线方向的运动性质与 a 相同，平行槽轴线方向做匀速直线运动。设滑块 b 的速度沿槽轴线和垂直槽轴线分速度为 $v_{bx}$、$v_{by}$，如图

\onepicture[h=4.2cm]{figs/17_an2.png}

当 $v_{bx}=0$ 时，滑块 b 第一次滑到最高点，由题意可知滑块 b 到达的最高点高度与滑块 a 的开始下滑的高度相等。此时速度为 $v_{by}$，经历的时间 $t_{1}$ 为

$t_{1}=\frac{1}{4}T=\frac{1}{4}\cdot2\pi\sqrt{\frac{r}{g}}=\frac{\pi}{2}\sqrt{\frac{r}{g}}$

又因与滑块 a 最高点相同，由题意可知 $v_{bx}=v_{by}=\frac{\sqrt{2}}{2}v_{b}=v_{0}$

即 $v_{b}=\sqrt{2}v_{0}=2\sqrt{gh}$

因为滑块 a、b 在最高点发生碰撞，设碰后滑块 b 的速度为 $v_{by}'$。

由动量守恒和机械能守恒有 $m_{b}v_{by}=m_{b}v_{by}'+m_{a}v_{a}$

$\frac{1}{2}m_{b}v_{by}^{2}=\frac{1}{2}m_{b}v_{by}'^{2}+\frac{1}{2}m_{a}v_{a}^{2}$

联立解得 $v_{by}'=-\sqrt{\frac{gh}{2}}$，$v_{a}=\sqrt{\frac{gh}{2}}$

（4）碰后滑块 a 在平行于槽轴线方向的速度始终为 $v_{a}$，从 $MN$ 边界射出的最基本的几种临界情况如图 1、2、3、4 所示。考虑周期性，则 $L$ 有多种情况

\onepicture[w=13cm]{figs/17_an3.png}

由题给数据可得 $t_{1}=\frac{\pi}{2}$，$v_{by}'=-0.4\Ums$，$v_{a}=0.4\Ums$

滑块 a 每一次在圆弧型斜面上滑或下滑的时间为 $t_{1}=\frac{\pi}{2}$

滑块 a 每一次滑过水平面的时间为 $t_{2}=\frac{d}{v_{0}}=\frac{0.4}{0.8}\Us=0.5\Us$

又由于滑块 a 从任一点出发回到该点同高度位置时的时间相等，设时间为 $T$，则 $T=4t_{1}+2t_{2}=(2\pi+1)\Us$

滑块 a 由碰后到从 $MN$ 之间飞出的时间 $t$ 满足 $t=\frac{L}{v_{a}}$，$L=0.4t$

所以由图 1、2 可知 $t_{1}<t<t_{1}+t_{2}$ 或 $t_{1}+T<t<t_{1}+t_{2}+T$，……

由图 3、4 可知 $3t_{1}+t_{2}<t<3t_{1}+2t_{2}$ 即 $t_{1}+\frac{T}{2}<t<t_{1}+t_{2}+\frac{T}{2}$

或者 $t_{1}+\frac{T}{2}+T<t<t_{1}+t_{2}+\frac{T}{2}+T$，……

综上分析可知 $t$ 应满足 $t_{1}+n\frac{T}{2}<t<t_{1}+t_{2}+n\frac{T}{2}$

则滑块 a 能从 $MN$ 之间飞出时 $L$ 的范围为

$[0.2\pi + n(0.4\pi + 0.2)]\Um < L < [0.2\pi + 0.2 + n(0.4\pi + 0.2)]\Um$，（$n = 0,1,2,3,4\cdots\cdots$）
}

\item
%% number: 18
%% paperName: 2026年1月浙江选考第 18 题 13 分
%% typeId: 6
%% type: 计算题
%% chapter: 磁场
%% point: 带电粒子在磁场中的运动
%% method: 无
%% score: 13
%% degree: 350
%% duplicateId: 0
%% body:
俄歇电子能谱（AES）广泛应用于材料表面成分分析。如图~\ref{20261月浙江18a} 所示，一束高能电子入射到样品表面，将某原子内层（如 K 层）的一个电子击出，形成一个空穴。随后，较外层（如 L 层）的一个电子跃迁至该空穴，并释放出能量，该能量可能以 X 光子的形式射出，也可能立即将另一核外电子（如 L 层或 M 层的电子）电离而逸出样品表面，该电子称为俄歇电子；现用电子动能 $E_{e} = 15.0\UkeV$ 的电子束轰击某样品表面，成功激发 KLM 俄歇过程（即初始空穴为 K 层、跃迁电子来自 L 层、逸出电子来自 M 层）和 KLL 俄歇过程（逸出电子来自 L 层）。已知该原子 K 层的电离能 $E_{K} \approx 10.0\UkeV$，L 层的电离能 $E_{L} \approx 1.0\UkeV$。已知电子的电荷量 $e = 1.6\times10^{-19}\UC$，电子质量 $m_{e} = 9.1\times10^{-31}\Ukg$，光速 $c = 3\times10^{8}\Ums$，普朗克常量 $h = 6.63\times10^{-34}\UJ\cdot\Us$。（计算结果保留一位有效数字）请回答：
\begin{enumerate}
	\item
	入射电子的德布罗意波长 $\lambda_{e}$。
	\item
	求射出的 X 光子的波长 $\lambda_{X}$；
	\item
	甲同学利用带电粒子在磁场中的运动规律，设计了如图~\ref{20261月浙江18b} 所示的测量俄歇电子动能的方案；俄歇电子从原点 $O$ 垂直 $y$ 轴和磁场方向进入匀强磁场，则 $y_{1} = 10.0\Ucm$ 和 $y_{2} = 10.5\Ucm$ 处被探测到，通过测得的俄歇电子的动能，求原子 M 层的电离能 $E_{M}$；
	\item
	乙同学认为用带电粒子在电场中的运动规律，测出俄歇电子的动能，请你帮乙同学设计一个方案，列出所需要测量的物理量，并给出计算俄歇电子动能的表达式。
\end{enumerate}
\twopicture[labelA=20261月浙江18a, labelB=20261月浙江18b, h=3.6cm]
{figs/18a.png}
{figs/18b.png}
%% answer: 
\jdanswer{
\begin{enumerate}
	\item
	$\lambda_{e} = 1\times10^{-11}\Um$

	\item
	$\lambda_{X} = 1\times10^{-10}\Um$

	\item
	$E_{M} \approx 0.2\UkeV$

	\item
	让待测电子束以水平初速度射入两块平行金属板之间，金属板长为 $L$，板间距为 $d$，两板之间加恒定的电压 $U$，形成匀强电场，在离开金属板右侧距离 $L_{1}$ 处垂直电子入射方向放置一荧光屏；需要测量的物理量：金属板长度 $L$，板间电压 $U$，板间距 $d$ 以及 $L_{1}$，电子飞出电场后打到荧光屏时垂直电子入射方向的侧移量 $y$；$E_{k} = \frac{eUL^{2}(L + 2L_{1})}{4dyL}$

\end{enumerate}
}
%% memo:
\memoanswer{
（1）由德布罗意公式 $\lambda_{e}=\frac{h}{p}$

又 $E_{k}=\frac{p^{2}}{2m}$

解得 $\lambda_{e}=\frac{h}{\sqrt{2m_{e}E_{k}}}=1\times10^{-11}\Um$

（2）因 $K_{\alpha}$ 线是由电子从 L 层跃迁到 K 层时释放的光子，则光子的能量为
$E = E_{K}-E_{L}=9\UkeV$

则由 $E = h\nu = \frac{hc}{\lambda}$

解得 $\lambda_{X}=\frac{hc}{E}=1\times10^{-10}\Um$

（3）电子在磁场中做圆周运动，洛伦兹力提供向心力 $evB = m\frac{v^{2}}{R}$

俄歇电子的动能为 $E_{k}=\frac{1}{2}mv^{2}=\frac{e^{2}B^{2}R^{2}}{2m}=\frac{e^{2}B^{2}}{8m}d^{2}$（$d$ 表示直径）

设 KLL 俄歇电子的动能为 $E_{KLL}$；KLM 俄歇电子的动能为 $E_{KLM}$，则有 $\frac{E_{KLL}}{E_{KLM}} = \frac{d_{1}^{2}}{d_{2}^{2}}$

KLL 俄歇电子：K 层空穴由 L 层电子填充，释放的能量为 $E_{K}-E_{L}$，并传递给 L 层另一个电子使其逸出，又消耗 $E_{L}$，则 KLL 俄歇电子的动能为 $E_{KLL} = (E_{K} - E_{L}) - E_{L} = 8.0\UkeV$

同理可得，KLM 俄歇电子的动能为 $E_{KLM} = (E_{K} - E_{L}) - E_{M}$

解得 $E_{M}=0.18\UkeV \approx 0.2\UkeV$

（4）如图，让待测电子束以水平初速度射入两块平行金属板之间，金属板长为 $L$，板间距为 $d$，两板之间加恒定的电压 $U$，形成匀强电场，在离开金属板右侧距离 $L_{1}$ 处垂直电子入射方向放置一荧光屏；

\onepicture[w=8.5cm]{figs/18_an.png}

需要测量的物理量：金属板长度 $L$，板间电压 $U$，板间距 $d$ 以及 $L_{1}$，电子飞出电场后打到荧光屏时垂直电子入射方向的侧移量 $y$；

计算方法：电子在水平方向做匀速直线运动，则 $t = \frac{L}{v}$

电子在竖直方向做初速度为零的匀加速直线运动，则加速度 $a=\frac{eU}{md}$

侧移量满足 $\frac{\frac{1}{2}at^{2}}{y}=\frac{\frac{1}{2}L}{\frac{1}{2}L+L_{1}}$

俄歇电子动能表达式 $E_{\mathrm{k}}=\frac{1}{2}mv^{2}=\frac{eUL^{2}(L+2L_{1})}{4dyL}$
}

\end{enumerate}

\end{document}
